\section{Exact local arithmetic and the two permutation carriers}
\label{sec:local-arithmetic}

\subsection{The surface envelope and field support}

An exact Macaulay resultant calculation for the four partial derivatives of
$F$, divided by the standard extraneous factor, gives
\begin{equation}
\begin{split}
\Delta_Y={}&2^{64}\cdot3^{43}\cdot5^7\cdot181^{24}\cdot283
\cdot997^{24}\cdot1801\\
&\hspace{18mm}\cdot2346241^{24}\cdot
14932047182473291995860108491583652133938007263719.
\end{split}
\label{eq:surface-divided-discriminant}
\end{equation}
The determinant was evaluated by two exact engines, one using fraction-free
elimination and one using exact polynomial arithmetic.  Their complete
matrices and determinant guards are summarized in \cref{app:local}.  The
factorization in \cref{eq:surface-divided-discriminant} is a bad-reduction
envelope for the integral surface model; it is not a field discriminant.

The global maximal order $\cO_E$ has rank $27$.  Its exact discriminant has
valuations
\begin{equation}
\begin{array}{c|rrrrrrrrr}
p&2&3&5&181&283&997&1801&2346241&q\\ \hline
v_p(\Disc E)&0&46&36&18&6&18&6&18&6.
\end{array}
\label{eq:discE-valuation-vector}
\end{equation}

\begin{proposition}[Envelope and exact support]
\label{prop:support}
The sets $\cB_Y$ and $\cR_{E,K}$ in \cref{thm:main} are respectively the
nine-prime set supported by \cref{eq:surface-divided-discriminant} and the
eight-prime set obtained by deleting $2$.
\end{proposition}

\begin{proof}
Away from the support of the divided discriminant, the integral cubic has
smooth reduction.  The finite line scheme then extends as a finite étale
scheme, so the Galois action on its geometric points is unramified.
\Cref{eq:discE-valuation-vector} shows that all eight displayed odd
primes ramify in $E$, while $2$ does not.

It remains to pass from the non-Galois field $E$ to its normal closure.  At
$2$, the zero permutation conductor says that inertia fixes all $27$ cosets
of a line stabilizer $H\leq\WE$.  The kernel of this action is
$\Core_{\WE}(H)$.  Since the line action in
\cref{eq:frozen-galois-input} is faithful, this core is trivial.  Thus
inertia at $2$ is trivial in $K$.  The same faithfulness shows that
nontrivial inertia at any of the eight odd primes cannot act trivially on all
$27$ lines.  Hence $E$ and $K$ have the same stated support.
\end{proof}

\subsection{Local maximal-order rows}

For a prime $\mathfrak p\mid p$ of $E$, write
\[
 e=e(\mathfrak p/p),\qquad f=f(\mathfrak p/p),\qquad
 d=v_{\mathfrak p}(\mathfrak D_{E_{\mathfrak p}/\Q_p}).
\]
Its contribution to $v_p(\Disc E)$ is $fd$.  Exact prime-ideal decomposition
in the maximal order gives the following complete rows.

\begin{table}[ht]
\centering
\caption{Maximal-order local rows.  Every row in the last line occurs at
each of the three rational primes, rather than being assigned one row per
prime.}
\label{tab:local-rows}
\begin{tabular}{@{}c l c c@{}}
\toprule
$p$&multiset of $(e,f,d)$&$\sum fd$&local degrees $ef$\\
\midrule
$3$&$\efrow{3}{1}{3},\efrow{6}{1}{7},\efrow{9}{1}{18}^2$&
$46$&$(3,6,9,9)$\\
$5$&$\efrow{1}{1}{0}^2,\efrow{5}{1}{7}^3,\efrow{10}{1}{15}$&
$36$&$(1,1,5,5,5,10)$\\
$181,997,2346241$&
$\efrow{3}{1}{2},\efrow{3}{2}{2},\efrow{3}{6}{2}$&
$18$&$(3,6,18)$\\
\bottomrule
\end{tabular}
\end{table}

The equalities
\[
3+7+18+18=46,\qquad 0+0+7+7+7+15=36,\qquad
2+2\cdot2+6\cdot2=18
\]
are direct different calculations, made before the representation-theoretic
conductor computation and used later as checks.  No constituent conductor is
obtained by splitting a global sum.

At the reflection primes $283,1801,q$, we deliberately do not list an
$(e,f,d)$ decomposition.  Their inertia is obtained by the geometric
ordinary-double-point argument in \cref{sec:reflection}.  The six
transpositions of a root reflection then give $v_p(\Disc E)=6$,
consistently with \cref{eq:discE-valuation-vector}.

\subsection{The 36-point authority and its boundary}

Let $\theta_{36}\in\Z[T]$ be the frozen degree-$36$ polynomial whose roots
label the double-sixes.  A second frozen polynomial $\delta_{36}$ has the
same global splitting field but is not a local authority in this paper.  The
distinction matters because a stable-looking finite-precision factorization
does not by itself identify the factors over $\Q_p$.

The factor-stability step uses the Krasner--Hensel criterion in the form
localized in \citet[Chapter II, \S2, Exercises 1--2, p.~30]
{serre1979local}.  For each asserted factorization, the working precision
exceeds both the valuation of the global polynomial discriminant and twice
the largest factor-polynomial discriminant valuation; the monic factors are
simple and multiply back to the original polynomial at the same precision.
The source supplies the criterion, while all exponents, factors, and
inequalities below are computations for \cref{eq:frozen-cubic}.

\begin{proposition}[Certified double-six partitions]
\label{prop:theta-authority}
The polynomial $\theta_{36}$ is a certified local authority with the records
in \cref{tab:theta-authority}.  No theorem assertion depends on
$\delta_{36}$, and no apparent $\delta_{36}$ factorization is used as
corroboration.
\end{proposition}

\begin{table}[ht]
\centering
\caption{Certified $\theta_{36}$ factor degrees and strict precision bounds.
The two bounds in the fourth column are the global
polynomial-discriminant valuation and twice the largest factor-polynomial
discriminant valuation.}
\label{tab:theta-authority}
\small
\begin{tabularx}{\textwidth}{@{}c c c c >{\raggedright\arraybackslash}X@{}}
\toprule
$p$&certified precisions&factor degrees&bounds&conclusion\\
\midrule
$3$&$(900,950,1000)$&$(3,3,3,9,18)$&$(886,538)$&
all three precisions clear both\\
$5$&$(900,950,1000)$&$(1,5,10,10,10)$&$(746,246)$&
all three precisions clear both\\
$181,997,2346241$&$(20,30,40)$&$(3,6,9,18)$&$(24,24)$&
precision $40$ clears both\\
\bottomrule
\end{tabularx}
\end{table}

For comparison only, at the tame $C_3$ primes the global
polynomial-discriminant valuation of $\delta_{36}$ is $840$, and twice its
largest factor-polynomial discriminant valuation is $408$.  Precision $40$
is below both values.  We therefore assign $\delta_{36}$ the formal role
\emph{bounded nonresult and nondependency}.  It supplies neither a premise
nor a cross-check, at the tame or wild primes.  Polynomial discriminants
appear nowhere else in the proof: they are separation bounds, not
substitutes for the field discriminants of $E$ or $K$.

\subsection{Why two actions are needed}

For $D\leq\WE$, its orbits on the $27$ lines have sizes equal to the local
factor degrees of $g$.  Its orbits on the $36$ double-sixes similarly match
the certified factor degrees of $\theta_{36}$.  Thus the wild targets are
\begin{equation}
\begin{array}{c|cc}
&27\text{-line carrier}&36\text{-double-six carrier}\\ \hline
p=3&(3,6,9,9)&(3,3,3,9,18)\\
p=5&(1,1,5,5,5,10)&(1,5,10,10,10).
\end{array}
\label{eq:wild-carrier-targets}
\end{equation}
The line carrier alone leaves multiple subgroup embeddings.  The double-six
carrier removes most of this ambiguity, but at $3$ three raw decomposition
classes still have the same simultaneous orbit targets.  Normality,
tame-quotient structure, branchwise differents, and the graded tame action
are therefore genuine proof steps, not database decoration.

The tame order-three case illustrates the same asymmetry.  Two $C_3$ classes
have compatible line behavior, but on double-sixes they have types
$3^{12}$ and $1^3 3^{11}$.  The certified $\theta_{36}$ type $3^{12}$
selects \Tom{6}; see \cref{prop:tame-c3}.  The $27$-line action then supplies
the $(V_6,V_{20})$ fixed spaces used for conductors.
